\subsection{Punto de partida y evolución del diseño}
\label{sec:proceso-diseno}

Antes de programar, la referencia de análisis y diseño fue la propuesta aprobada. No se definieron requisitos ni decisiones adicionales a los recogidos en ella, ni se elaboraron diagramas UML previos. Se planteaba identificar indicios de phishing en correos o URLs mediante contenido, enlaces y cabeceras, aplicar heurísticas, integrar fuentes externas y presentar resultados comprensibles. El resumen añadía certificados y recomendaciones preventivas. Se preveían Python, procesamiento de correo y HTML, Streamlit y bibliotecas de aprendizaje automático; la arquitectura final quedó por concretar.

Durante la implementación, ese planteamiento se transformó en decisiones concretas. La primera versión registrada separa una interfaz Streamlit del analizador y las heurísticas. Esta organización permitía recoger el mensaje, examinar indicios y mostrar una valoración sin concentrar todas las funciones en la interfaz. La incorporación posterior del clasificador TF-IDF + MLP amplió el análisis basado en reglas y dio lugar a los modos heurístico, neuronal y combinado. Su selección de biblioteca y sus límites se explican en el apartado~\ref{sec:decision-tensorflow}; no fueron una elección cerrada en la propuesta.

La incorporación de Gmail, Telegram y la extensión amplió las formas de entrada y notificación. El diseño evolucionó hacia un backend HTTP común para concentrar el análisis y mantener una versión activa de cada modelo compartida por los clientes. También se concretó la limitación al análisis estático: quedaron fuera la reputación externa, la validación de certificados y la inspección activa de sitios web. El apartado~\ref{sec:objetivos-propuesta} recoge las razones y las consecuencias de esas exclusiones. Las iteraciones de evaluación añadieron criterios de reproducibilidad y calibración, junto con las comprobaciones descritas en los capítulos 4 y 6.

La tabla~\ref{tab:5-1} expresa el resultado de esta evolución: RF-01 y RF-03 concretan el análisis y la presentación previstos, aunque Gmail se incorporó durante el desarrollo; RF-02, RF-04 y RF-05 describen los modos, la centralización y las alertas del prototipo final. Los RNF formalizan las condiciones adoptadas para verificarlo. Los cinco modelos UML del apartado~\ref{sec:uml} documentan cómo se relacionan esas funciones, clases, interacciones, componentes y procesos en la versión entregada. Esta descripción retrospectiva complementa la metodología del apartado~\ref{sec:metodologia} y permite explicar el proceso seguido sin atribuir a la fase inicial documentación elaborada después.
